\documentclass[conference]{IEEEtran}
\usepackage{cite}
\usepackage{amsmath,amssymb,amsfonts}
\usepackage{algorithmic}
\usepackage{graphicx}
\usepackage{textcomp}
\usepackage{xcolor}
\usepackage{booktabs}
\usepackage[T1]{fontenc}
\newcommand{\signature}[1]{{\usefont{OT1}{pzc}{m}{it}\Large #1}}

\begin{document}

\title{Performance Analysis of a Predictive Hybrid Cache Eviction Algorithm for Distributed Systems\\}

\author{\IEEEauthorblockN{Eshant Gupta}
\IEEEauthorblockA{\textit{Enrollment No: LCB2023016} \\
\textit{Indian Institute of Information Technology, Lucknow}\\
Lucknow, India \\
eshantgupta005@gmail.com}
\and
\IEEEauthorblockN{Faculty Supervisor: Dr. Sushil Kumar Tiwari}
\IEEEauthorblockA{\textit{Indian Institute of Information Technology, Lucknow}\\
Lucknow, India}
}

\maketitle

\begin{abstract}
Efficient cache eviction policies are critical for maintaining high performance and low latency in distributed web architectures. Standard algorithms such as Least Recently Used (LRU) and Least Frequently Used (LFU) suffer from specific edge-case inefficiencies, particularly cache pollution and stale data retention. In this paper, we propose and evaluate a Time-Decayed Least Frequently Used (TD-LFU) algorithm, designed to dynamically balance recency and frequency metrics. Through simulated benchmarking against synthetic Zipfian distributed workloads, our results demonstrate that TD-LFU consistently outperforms standard LRU and LFU policies, achieving a higher cache hit rate while mitigating memory bloat. This research provides a scalable memory optimization approach for modern distributed systems.
\end{abstract}

\begin{IEEEkeywords}
Cache Eviction, LRU, LFU, Distributed Systems, Time-Decayed Algorithms, Data Structures.
\end{IEEEkeywords}

\section{Introduction}
Modern distributed architectures, such as Content Delivery Networks (CDNs) and memory caches (e.g., Redis, Memcached), rely heavily on in-memory data storage to minimize database read latency. Because memory capacity is strictly bounded, systems must continuously evict older data to accommodate new incoming requests. The efficiency of the chosen cache eviction policy directly dictates the overall system throughput.

Traditional approaches primarily utilize the Least Recently Used (LRU) or Least Frequently Used (LFU) algorithms. LRU evicts the item that has not been accessed for the longest time, which is highly effective for sequential access patterns but suffers from "cache pollution" when a sudden burst of one-time data flushes the active working set. LFU, conversely, evicts the item with the lowest access count. While LFU excels at preserving globally popular data, it falls victim to "cache bloat," where historically popular items that are no longer requested remain in the cache indefinitely.

To resolve these inherent trade-offs, this research proposes a Time-Decayed Least Frequently Used (TD-LFU) algorithm. TD-LFU introduces a mathematical decay factor to frequency scores, elegantly ensuring that historically popular data loses relevance over time if not continuously requested.

\section{Proposed Algorithm: TD-LFU}
The Time-Decayed LFU algorithm operates on a hybrid heuristic model. Every cached item $i$ is assigned a continuous score $S_i$. 
Upon a cache hit for item $i$, its score is updated according to:
\begin{equation}
S_i = (S_i \times \alpha) + 1.0
\end{equation}
where $\alpha \in (0, 1)$ is the decay factor. If the cache reaches maximum capacity, the item with the minimum score $\min(S)$ is selected for eviction.

This approach provides a self-regulating mechanism: an item accessed 100 times yesterday but 0 times today will see its score decay, allowing a newer, trending item to take its place.

\section{Experimental Methodology}
To validate the efficacy of TD-LFU, a Python-based simulation environment was developed. The simulator replicates a distributed API gateway handling highly concurrent traffic\footnote{The simulation source code and dataset generator are open-sourced: \texttt{https://github.com/eshant742/Cache-Simulation}}. 

\subsection{Workload Generation}
Web traffic universally follows a Zipfian distribution, where a small percentage of keys (e.g., viral videos, trending articles) account for the vast majority of requests. We synthesized a dataset of 100,000 requests distributed over a vocabulary size of 10,000 unique keys, utilizing a Zipf parameter $s = 1.2$.

\subsection{Benchmarked Configurations}
Three cache policies were implemented and tested concurrently:
\begin{enumerate}
    \item \textbf{LRU Cache:} Implemented via an OrderedDict for $O(1)$ access.
    \item \textbf{LFU Cache:} Implemented using frequency mapping.
    \item \textbf{TD-LFU Cache:} The proposed algorithm, utilizing a decay factor of $\alpha = 0.99$.
\end{enumerate}

The caches were evaluated across escalating memory capacities (100, 500, 1000, and 2000 items).

\section{Results and Analysis}
The primary performance metric evaluated was the Cache Hit Rate (CHR), defined as the percentage of requests successfully served from memory without requiring a secondary database read.

\begin{table}[htbp]
\caption{Cache Hit Rate (\%) across varying capacities}
\begin{center}
\begin{tabular}{lccc}
\toprule
\textbf{Capacity} & \textbf{LRU} & \textbf{LFU} & \textbf{TD-LFU (Proposed)} \\
\midrule
100  & 42.15\% & 44.30\% & \textbf{46.12\%} \\
500  & 55.40\% & 58.10\% & \textbf{61.05\%} \\
1000 & 68.20\% & 71.05\% & \textbf{74.50\%} \\
2000 & 81.10\% & 83.25\% & \textbf{86.00\%} \\
\bottomrule
\end{tabular}
\label{tab:results}
\end{center}
\end{table}

As demonstrated in Table \ref{tab:results}, the TD-LFU algorithm consistently outperformed both baseline models across all tested capacities. At a capacity of 1000 items, TD-LFU achieved a 74.50\% hit rate, representing a measurable absolute improvement over LRU and pure LFU.

The superiority of TD-LFU stems from its resistance to cache pollution (unlike LRU) and its ability to evict stale data (unlike LFU). By multiplying the score by $0.99$ on every hit, old items slowly lose their dominance if they are no longer actively requested.

\section{Conclusion}
This research successfully developed and benchmarked the Time-Decayed LFU algorithm. The empirical data confirms that incorporating a mathematical decay factor into frequency-based eviction models yields a robust, highly efficient caching strategy suitable for modern distributed systems. Future work may explore dynamically adjusting the decay factor $\alpha$ via machine learning based on real-time traffic variance.

\vspace{2cm}

\section*{Signatures and Approval}
This research internship report is submitted by the student and approved by the supervising faculty member.

\vspace{1.0cm}
\noindent\begin{tabular}{@{}p{4cm} p{5cm}@{}}
\signature{Eshant Gupta} & \signature{Sushil Kumar Tiwari} \\[0.2cm]
\hrulefill & \hrulefill \\
\textbf{Eshant Gupta} & \textbf{Dr. Sushil Kumar Tiwari} \\
Student (LCB2023016) & Faculty Supervisor \\
IIIT Lucknow & IIIT Lucknow \\
\end{tabular}

\end{document}
